% This presentation is based on a Beamer theme from Seth Brown, distributed
% under the following license:
%
% ----------------------------------------------------------------------------
% This program can be redistributed and/or modified under the terms
% of the GNU Public License, version 3.
%
% Seth Brown, Ph.D.
% sethbrown@drbunsen.org

\documentclass[t]{beamer}

\usepackage{hyperref} % urls
\usepackage{graphicx} % graphics
\usepackage[utf8x]{inputenc}
\usepackage{fancyvrb}
\usepackage{tcolorbox}
\usepackage{animate}

% suppress navigation bar
\beamertemplatenavigationsymbolsempty

\mode<presentation>
{
  \usetheme{codeplay}
  \setbeamercovered{transparent}
  %\setbeamertemplate{itemize item}[circle]
  \setbeamertemplate{itemize item}{\small\raise0.8pt\hbox{$\bullet$}}
  %\setbeamertemplate{itemize subitem}{\tiny\raise1.5pt\hbox{\donotcoloroutermaths$\blacktriangleright$}}
  \setbeamertemplate{itemize subitem}{\tiny\raise1.5pt\hbox{$\blacktriangleright$}}
  \setbeamertemplate{section in toc shaded}[default][100]
  \setbeamertemplate{subsection in toc shaded}[default][100]

\setbeamertemplate{subsection in toc}{%
\leavevmode\leftskip=4.00ex%
  \llap{\raisebox{0.15ex}{\textcolor{structure}{\small$\bullet$}}\kern1ex}%
  \inserttocsubsection\par%
}
}

% set fonts
\usepackage{fontspec}
\setsansfont{Calibri}
\setmonofont{Consolas}
\setbeamerfont{frametitle}{size=\LARGE,series=\bfseries}

% color definitions
\usepackage{color}
\definecolor{uiwhite}{RGB}{255, 255, 255}
\definecolor{uidarkgray}{RGB}{24,24,25}
\definecolor{uigray}{RGB}{51,51,51}
\definecolor{uilightgray}{RGB}{123,123,123}
\definecolor{uicyan}{RGB}{0,255,255}
\definecolor{uipink}{RGB}{255,153,255}
\definecolor{uipalegreen}{RGB}{153,255,153}
\definecolor{uiorange}{RGB}{237,175,28}
\hypersetup{colorlinks,urlcolor=uicyan,linkcolor=uiwhite}

% title slide definition
\title{Creating an SPMD Vectorizer for OpenCL with LLVM}
\subtitle{LLVM 2015 tutorial}
\author{Pierre-André Saulais}
\institute{Codeplay Software \\ @codeplaysoft}

\date{\today}

% Code box formatting
\fvset{fontsize=\scriptsize,tabsize=4}
\tcbset{colback=uigray,colframe=uilightgray,colupper=uiwhite,
        left=1pt,right=1pt,top=1pt,bottom=1pt,arc=0pt,
        toprule=1pt,bottomrule=1pt,leftrule=1pt,rightrule=1pt}

\newenvironment{codebox}
 {\VerbatimEnvironment
  \begin{tcolorbox}%
  \begin{Verbatim}}
 {\end{Verbatim}\end{tcolorbox}}

% Code caption that comes after the code box.
\newcommand{\codecaption}[2][-4.4ex]{%
  \vspace{#1}
  \hfill\mbox{\footnotesize{\textcolor{uicyan}{#2}}\hspace{0.5em}}
  \vspace*{2ex}}

\newcommand{\sourcebox}[3][firstline=1]{%
  \begin{tcolorbox}\VerbatimInput[#1]{#2}\end{tcolorbox}
  \codecaption{#3}}

\newcommand{\examplebox}[2][firstline=1]{\sourcebox[#1]{examples/#2}{#2}}

\newcommand{\codeempha}[1]{\textcolor{uipink}{#1}}
\newcommand{\codeemphb}[1]{\textcolor{uicyan}{#1}}
\newcommand{\codeemphc}[1]{\textcolor{uipalegreen}{#1}}
\newcommand{\codeemphd}[1]{\textcolor{uiorange}{#1}}

\newcommand{\varying}[1]{\codeempha{#1}}
\newcommand{\uniform}[1]{\codeemphb{#1}}
\newcommand{\idx}[1]{[#1]}

%\newcommand{\penguin}[0]{\includegraphics[width = 1em]{images/penguin.png}}

\newcommand{\penguin}[0]{\animategraphics[loop,autoplay, width = 1em]{12}{images/penguin-}{0}{19}\hspace{0.1em}}

% Create a slide for a part.
\newcommand{\talkpart}[2]{%
\section{Part #1: #2}

\begin{frame}[c]{Part #1}

\vspace{1cm}
\centerline{\LARGE{#2}}

\end{frame}}

% Create a slide for a part section.
\newcommand{\talksection}[1]{%
\subsection{#1}
\begin{frame}{\insertsectionhead}
\vspace{2ex}
\begin{NoHyper}
\tableofcontents[sectionstyle=hide,subsectionstyle=show/shaded/hide]
\end{NoHyper}
\end{frame}}

%%%%%%%%%%%%%%%%%%%%%%%%%%%%%%%%%%%%%%%%%%%%%%%%%%%%%%%%%%%%%%%%%%%%%%%%%%%%%%%%

\begin{document}

\setbeamertemplate{background}
{\includegraphics[width=\paperwidth,height=\paperheight]{dark_background_title.png}}
\setbeamertemplate{footline}[default]

\begin{frame}
  \vspace{4ex}
  \titlepage
\end{frame}

%%%%%%%%%%%%%%%%%%%%%%%%%%%%%%%%%%%%%%%%%%%%%%%%%%%%%%%%%%%%%%%%%%%%%%%%%%%%%%%%

% Set the background for the rest of the slides.
%\setbeamertemplate{background}
% {\includegraphics[width=\paperwidth,height=\paperheight]{dark_background.png}}
\setbeamertemplate{background}{}
\setbeamertemplate{footline}[codeplaytheme]

\section*{Introduction}

\begin{frame}{What this tutorial is about}

\begin{itemize}
    \item Vectorizing
    \begin{itemize}
        \item Transform whole functions using LLVM
        \item "Horizontal" vectorization
    \end{itemize}  
    \item SPMD programs
    \begin{itemize}
        \item Data-parallel execution model
        \item Used with compute frameworks like OpenCL and CUDA
    \end{itemize}
    \item for CPU-like processors, i.e. each core has
    \begin{itemize}
        \item A program counter
        \item Vector SIMD unit(s)
    \end{itemize}
\end{itemize}

\end{frame}

%%%%%%%%%%%%%%%%%%%%%%%%%%%%%%%%%%%%%%%%%%%%%%%%%%%%%%%%%%%%%%%%%%%%%%%%%%%%%%%%

\begin{frame}{Overview}
\tableofcontents
\end{frame}

%%%%%%%%%%%%%%%%%%%%%%%%%%%%%%%%%%%%%%%%%%%%%%%%%%%%%%%%%%%%%%%%%%%%%%%%%%%%%%%%

%\talkpart{1}{Background}
%%%%%%%%%%%%%%%%%%%%%%%%%%%%%%%%%%%%%%%%%%%%%%%%%%%%%%%%%%%%%%%%%%%%%%%%%%%%%%%%%

\begin{frame}{SPMD execution model}

\end{frame}


%% Kernel function has no return value, takes in buffers (arrays)
%% Vectorization does not change the signature for kernels
%% Can also vectorize "normal" functions, but this results in signature changes


\talkpart{2}{Implementing a SPMD Vectorizer}
%%%%%%%%%%%%%%%%%%%%%%%%%%%%%%%%%%%%%%%%%%%%%%%%%%%%%%%%%%%%%%%%%%%%%%%%%%%%%%%%

\talksection{Overview}

\begin{frame}{SPMD Vectorizer}

\begin{itemize}
    \item Is used to vectorize a SPMD program's entry point function
    \item Given a function $F$ and vectorization factor $N$, produces a function $VF_N$
    \begin{itemize}
        \item Calling $VF_N$ is like calling $F$, but $N$ times (on consecutive work-items)
        \item $F$ and $VF_N$ have the same signature
    \end{itemize}
    
    \item The original function can be preserved or not
    \begin{itemize}
        \item In-place: vectorize the original function, no need for cloning
        \item Transform cloned function: allow vectorization to fail
    \end{itemize}
    
    \item Vectorization may be allowed to fail
    \begin{itemize}
        \item On failure, the original function can be used 
    \end{itemize}
\end{itemize}

\end{frame}

%%%%%%%%%%%%%%%%%%%%%%%%%%%%%%%%%%%%%%%%%%%%%%%%%%%%%%%%%%%%%%%%%%%%%%%%%%%%%%%%

\begin{frame}{Structure}

\begin{itemize}
    \item Pipeline structure
    \begin{itemize}
        \item Function to vectorize is repeatedly transformed by different stages
        \item Stages are independent from each other
        \item Each stage consists of one or more IR passes
        \item Most stages require some analysis
    \end{itemize}
    
    \item The LLVM pass manager (both new and original) is a good fit here
    \begin{itemize}
        \item Most passes work at the function level
        \item Access to the module is sometimes needed
    \end{itemize}
\end{itemize}

\end{frame}

%%%%%%%%%%%%%%%%%%%%%%%%%%%%%%%%%%%%%%%%%%%%%%%%%%%%%%%%%%%%%%%%%%%%%%%%%%%%%%%%

\begin{frame}[c]{Stages}

[Stage diagram with stages:\\ preparation, CFG conversion, scalarization, packetization, optimizations]

\end{frame}

%%%%%%%%%%%%%%%%%%%%%%%%%%%%%%%%%%%%%%%%%%%%%%%%%%%%%%%%%%%%%%%%%%%%%%%%%%%%%%%%

\begin{frame}{Analyses}

\begin{itemize}
    \item Capture information about the IR to vectorize
    \item May be invalidated after IR transformations
    \item May depend on other analyses
\end{itemize}

\begin{itemize}
    \item Uniform Value Analysis
    \begin{itemize}
        \item Marks values as either \uniform{uniform} or \varying{varying}
    \end{itemize}

    \item Control Flow Analysis
    \begin{itemize}
        \item Determines which basic blocks are divergent
        \item Builds a Control Dependency Graph
    \end{itemize}
\end{itemize}

\end{frame}

%%%%%%%%%%%%%%%%%%%%%%%%%%%%%%%%%%%%%%%%%%%%%%%%%%%%%%%%%%%%%%%%%%%%%%%%%%%%%%%%

\begin{frame}{Implementation strategy}

\begin{itemize}
    \item Create test kernels
    \begin{itemize}
        \item Start with very simple kernels (e.g. copy buffer, add two buffers)
        \item Gradually add more features (e.g. non-sequential memory accesses, vector instructions, etc)
    \end{itemize}

    
    \item Suggested implementation order
    \begin{itemize}
        \item Preparation and packetization first (required for simplest kernels)
        \item Then easier features: builtins, memory addressing, scalarization, instantiation
        \item More complex features last: control flow, optimizations
    \end{itemize}
\end{itemize}

\end{frame}

%%%%%%%%%%%%%%%%%%%%%%%%%%%%%%%%%%%%%%%%%%%%%%%%%%%%%%%%%%%%%%%%%%%%%%%%%%%%%%%%

\talksection{Packetization Stage}

\begin{frame}{Packetization Overview}

\begin{itemize}
    \item Stage that does the actual vectorization
    \begin{itemize}
        \item Needs a vectorization factor $N$ (SIMD width)
        \item Calling $VF_N$ is like calling $F$, but $N$ times
        \item Straightforward thanks to preparation from previous stages
    \end{itemize}
    
    \item This is done per-instruction, for the whole function
    \begin{itemize}
        \item Instructions that define a value: define $N$ values, one for each instance
        \item Instructions with side effects: perform side effects for each instance
    \end{itemize}
    
    \item Only \varying{varying} instructions need packetization
    \begin{itemize}
        \item \uniform{Uniform} instructions can remain scalar, executed once per work-group
        \item Requires \emph{Uniform Value Analysis} to know which instructions to vectorize
    \end{itemize}
    
\end{itemize}

\end{frame}

%%%%%%%%%%%%%%%%%%%%%%%%%%%%%%%%%%%%%%%%%%%%%%%%%%%%%%%%%%%%%%%%%%%%%%%%%%%%%%%%

\begin{frame}{Uniform Value Analysis}

\begin{itemize}
    \item Finds 'root' values
    \begin{itemize}
        \item \varying{Varying} values with no \varying{varying} operand
        \item Example: \varying{get\_global\_id(0)} has a different value for each isntance
    \end{itemize}
    \item Marks each IR value as \uniform{uniform} or \varying{varying}
    \begin{itemize}
        \item All values start as \uniform{uniform}
        \item Marking a value as \varying{varying} causes all users to also be marked \varying{varying}
        \item Marking is done recursively, starting with roots
        \item Values are marked before their users, to avoid cycles (phi nodes)
    \end{itemize}
\end{itemize}

\end{frame}

%%%%%%%%%%%%%%%%%%%%%%%%%%%%%%%%%%%%%%%%%%%%%%%%%%%%%%%%%%%%%%%%%%%%%%%%%%%%%%%%

\begin{frame}[fragile]{Uniform Value Analysis}

Example that combines \uniform{uniform} and \varying{varying} values:

\begin{codebox}[commandchars=\\\[\]]
kernel void add_uniform(global int *\uniform[dst], global int *\uniform[src], int \uniform[alpha]) {
    int \varying[tid] = \varying[get_global_id](0);
    \uniform[dst]\idx[\varying[tid]] = \uniform[src]\idx[\varying[tid]] + (\uniform[alpha] - \uniform[1]);
}
\end{codebox}

\begin{codebox}[commandchars=\\\[\]]
define void @add_uniform(i32* \uniform[%dst], i32* \uniform[%src], i32 \uniform[%alpha]) {
entry:
  \varying[%tid] = i32 \varying[@get_global_id(i32 0)]
  \varying[%arrayidx] = getelementptr i32* \uniform[%src], i32 \varying[%tid]
  \varying[%tmp] = load i32* \varying[%arrayidx], align 4
  \uniform[%sub] = sub i32 \uniform[%alpha], \uniform[1]
  \varying[%add] = add i32 \uniform[%sub], \varying[%tmp]
  \varying[%arrayidx2] = getelementptr i32* \uniform[%dst], i32 \varying[%tid]
  store i32 \varying[%add1], i32* \varying[%arrayidx2], align 4
  ret void
}
\end{codebox}

\end{frame}

%%%%%%%%%%%%%%%%%%%%%%%%%%%%%%%%%%%%%%%%%%%%%%%%%%%%%%%%%%%%%%%%%%%%%%%%%%%%%%%%

\begin{frame}[c]{Uniform Value Analysis}

[IR graph to show values, highlight \uniform{uniform} and \varying{varying} nodes]

\end{frame}

%%%%%%%%%%%%%%%%%%%%%%%%%%%%%%%%%%%%%%%%%%%%%%%%%%%%%%%%%%%%%%%%%%%%%%%%%%%%%%%%

\begin{frame}{Packetization Process}

\begin{itemize}
    \item Find leaves
    \begin{itemize}
        \item Leaves allow \varying{varying} values to 'escape' from the function, they are:
        \item Store instructions (\varying{varying} operand)
        \item Call instructions (\varying{varying} operand, or call has no use)
        \item Return instructions
    \end{itemize}
\end{itemize}

\begin{itemize}
    \item Recursively packetize leaves and their operands
    \begin{itemize}
        \item Broadcast \uniform{uniform} values (e.g. argument, constants)
        \item Replace \varying{get\_global\_id(0)} with a vector of IDs
        \item Packetize operands first, then instruction (top-down)
        \item Cache packetized values to prevent duplication
    \end{itemize}
\end{itemize}

\begin{itemize}
    \item Delete original scalar instructions if dead
\end{itemize}

\end{frame}

%%%%%%%%%%%%%%%%%%%%%%%%%%%%%%%%%%%%%%%%%%%%%%%%%%%%%%%%%%%%%%%%%%%%%%%%%%%%%%%%

\begin{frame}[c]{Packetization Example}

[Packetization example \\
TODO: original graph \\
TODO: vectorized graph?]

\end{frame}

%%%%%%%%%%%%%%%%%%%%%%%%%%%%%%%%%%%%%%%%%%%%%%%%%%%%%%%%%%%%%%%%%%%%%%%%%%%%%%%%

\begin{frame}[fragile]{Packetization Example}

\begin{codebox}[commandchars=\\\[\]]
define void @__v4_add_uniform(i32* \uniform[%dst], i32* \uniform[%src], i32 \uniform[%alpha]) {
entry:
  \varying[%tid] = call i32 \varying[@get_global_id(i32 0)]
  \varying[%arrayidx] = getelementptr i32* \uniform[%src], i32 \varying[%tid]
  \varying[%0] = bitcast i32* \varying[%arrayidx] to <4 x i32>*
  \varying[%1] = load <4 x i32>* \varying[%0], align 4

  ; Broadcast (alpha - 1) to a vector
  \uniform[%sub] = sub i32 \uniform[%alpha], \uniform[1]
  \uniform[%insert] = insertelement <4 x i32> undef, i32 \uniform[%sub], i32 0
  \uniform[%broadcast_sub] = shufflevector <4 x i32> \uniform[%insert], ...

  \varying[%add] = add nsw <4 x i32> \uniform[%broadcast_sub], \varying[%1]

  \varying[%arrayidx2] = getelementptr i32* \uniform[%dst], i32 \varying[%tid]
  \varying[%2] = bitcast i32* \varying[%arrayidx2] to <4 x i32>*
  store <4 x i32> \varying[%add], <4 x i32>* \varying[%2], align 4
  ret void
}
\end{codebox}

\end{frame}

%%%%%%%%%%%%%%%%%%%%%%%%%%%%%%%%%%%%%%%%%%%%%%%%%%%%%%%%%%%%%%%%%%%%%%%%%%%%%%%%

\begin{frame}[fragile]{Packetizing Phi nodes}

\begin{itemize}
    \item Phi nodes can introduce cycles in the use/def graph
    \begin{itemize}
        \item Phi operands that reference the same phi
        \item Example: reduction variables
    \end{itemize}
    \item Break the cycle by creating 'empty' phi nodes
    \item Packetize operands of phi nodes after leaves have been packetized
\end{itemize}

[TODO: graph]

\end{frame}

%%%%%%%%%%%%%%%%%%%%%%%%%%%%%%%%%%%%%%%%%%%%%%%%%%%%%%%%%%%%%%%%%%%%%%%%%%%%%%%%

\begin{frame}{Memory Addressing}

\begin{itemize}
    %\item \varying{get\_global\_id(0)} was not packetized to a sequence of IDs. Why?
    \item Loads and stores are packetized according to the addressing mode
    \begin{itemize}
        \item Each operation can access $N$ memory elements. How are they laid out?
        \item Different layouts result in different packetized operations
        \item Most common layouts can be described with a single stride
    \end{itemize}
    
    \item Stride is the distance between successive elements
    \begin{itemize}
        \item Expressed in number of elements
        \item "How many elements are skipped in memory to get to the next one"
        \item One means elements are consecutive
    \end{itemize}
    
    \item Not all kinds of memory operations can be expressed in IR
    \begin{itemize}
        \item Generate calls to internal builtins
        \item Internal builtins can be implemented for each target as supported 
    \end{itemize}
    
\end{itemize}

\end{frame}

%%%%%%%%%%%%%%%%%%%%%%%%%%%%%%%%%%%%%%%%%%%%%%%%%%%%%%%%%%%%%%%%%%%%%%%%%%%%%%%%

\begin{frame}{Memory Addressing Modes}

\begin{itemize}
    \item $Stride = 0$
    \begin{itemize}
        \item Offset is \uniform{uniform}
        \item Generate a regular \emph{scalar} load or store
    \end{itemize}
    
    \item $Stride = 1$
    \begin{itemize}
        \item Packetized offset is a sequence like $<0, 1, 2, 3>$
        \item Generate a regular \emph{vector} load or store
    \end{itemize}
    
    \item $Stride > 1$
    \begin{itemize}
        \item Packetized offset is a sequence like $<0, 4, 8, 12>$
        \item Generate an \emph{interleaved} load or store
    \end{itemize}
    
    \item Variable stride
    \begin{itemize}
        \item Packetized offset can be any sequence
        \item Generate a \emph{gather} load or \emph{scatter} store
    \end{itemize}
    
    \item Handle negative strides by reversing the vector
\end{itemize}

\end{frame}

%%%%%%%%%%%%%%%%%%%%%%%%%%%%%%%%%%%%%%%%%%%%%%%%%%%%%%%%%%%%%%%%%%%%%%%%%%%%%%%%

\talksection{Scalarization Stage}

\begin{frame}{Scalarization Overview}

\begin{itemize}
    \item What does it do?
    \item Requires scalarization analysis
\end{itemize}

\end{frame}

%%%%%%%%%%%%%%%%%%%%%%%%%%%%%%%%%%%%%%%%%%%%%%%%%%%%%%%%%%%%%%%%%%%%%%%%%%%%%%%%

\begin{frame}{Scalarization Analysis}

\begin{itemize}
    \item Looks for vector instructions
    \begin{itemize}
        \item Leaves that define vector values, vector stores
        \item Vector extractions
        \item Vector -> scalar bitcasts
    \end{itemize}
    
\end{itemize}

\end{frame}

%%%%%%%%%%%%%%%%%%%%%%%%%%%%%%%%%%%%%%%%%%%%%%%%%%%%%%%%%%%%%%%%%%%%%%%%%%%%%%%%

\begin{frame}{Scalarization Process}


\end{frame}

%%%%%%%%%%%%%%%%%%%%%%%%%%%%%%%%%%%%%%%%%%%%%%%%%%%%%%%%%%%%%%%%%%%%%%%%%%%%%%%%

\begin{frame}{Scalarization Example}


\end{frame}

%%%%%%%%%%%%%%%%%%%%%%%%%%%%%%%%%%%%%%%%%%%%%%%%%%%%%%%%%%%%%%%%%%%%%%%%%%%%%%%%

\talksection{Control Flow Conversion Stage}

\begin{frame}{Control Flow Conversion Overview}

\begin{itemize}
    \item What does it do?
    \item Why is it needed?
\end{itemize}

\end{frame}

%%%%%%%%%%%%%%%%%%%%%%%%%%%%%%%%%%%%%%%%%%%%%%%%%%%%%%%%%%%%%%%%%%%%%%%%%%%%%%%%

\begin{frame}{Control Flow Conversion: if}

\end{frame}

%%%%%%%%%%%%%%%%%%%%%%%%%%%%%%%%%%%%%%%%%%%%%%%%%%%%%%%%%%%%%%%%%%%%%%%%%%%%%%%%

\begin{frame}{Mask Generation}

\end{frame}

%%%%%%%%%%%%%%%%%%%%%%%%%%%%%%%%%%%%%%%%%%%%%%%%%%%%%%%%%%%%%%%%%%%%%%%%%%%%%%%%

\begin{frame}{Applying Masks}

\end{frame}

%%%%%%%%%%%%%%%%%%%%%%%%%%%%%%%%%%%%%%%%%%%%%%%%%%%%%%%%%%%%%%%%%%%%%%%%%%%%%%%%

\begin{frame}{Masked Memory Operations}

\end{frame}

%%%%%%%%%%%%%%%%%%%%%%%%%%%%%%%%%%%%%%%%%%%%%%%%%%%%%%%%%%%%%%%%%%%%%%%%%%%%%%%%

\begin{frame}{Phi Conversion}

\end{frame}

%%%%%%%%%%%%%%%%%%%%%%%%%%%%%%%%%%%%%%%%%%%%%%%%%%%%%%%%%%%%%%%%%%%%%%%%%%%%%%%%

\begin{frame}{CFG Linearization}

\end{frame}

%%%%%%%%%%%%%%%%%%%%%%%%%%%%%%%%%%%%%%%%%%%%%%%%%%%%%%%%%%%%%%%%%%%%%%%%%%%%%%%%

\begin{frame}{Control Flow Conversion: loops}

\end{frame}

%%%%%%%%%%%%%%%%%%%%%%%%%%%%%%%%%%%%%%%%%%%%%%%%%%%%%%%%%%%%%%%%%%%%%%%%%%%%%%%%

\begin{frame}{Finding Loop Live Variables}

\end{frame}

%%%%%%%%%%%%%%%%%%%%%%%%%%%%%%%%%%%%%%%%%%%%%%%%%%%%%%%%%%%%%%%%%%%%%%%%%%%%%%%%

\begin{frame}{Merging Loop Live Variables}

\end{frame}


%\talkpart{3}{Going further}
%%%%%%%%%%%%%%%%%%%%%%%%%%%%%%%%%%%%%%%%%%%%%%%%%%%%%%%%%%%%%%%%%%%%%%%%%%%%%%%%%

\begin{frame}{SIMD width detection}

\end{frame}

%%%%%%%%%%%%%%%%%%%%%%%%%%%%%%%%%%%%%%%%%%%%%%%%%%%%%%%%%%%%%%%%%%%%%%%%%%%%%%%%

\begin{frame}{Vectorizing builtin function calls}

% Builtin: the vectorizer has some knowledge of scalar -> vector function mapping

\end{frame}

%%%%%%%%%%%%%%%%%%%%%%%%%%%%%%%%%%%%%%%%%%%%%%%%%%%%%%%%%%%%%%%%%%%%%%%%%%%%%%%%

\begin{frame}{Vectorizing builtin functions}

% By vectorizing the builtin's body
% This changes the function signature (return value, some arguments)
% Builtin: the vectorizer has some knowledge of which arguments need packetization
% Need argument placeholders (cloning required)
% Packetized arguments are roots (Uniform Value Analysis)
% Return instructions are leaves (Packetization Stage)

\end{frame}

%%%%%%%%%%%%%%%%%%%%%%%%%%%%%%%%%%%%%%%%%%%%%%%%%%%%%%%%%%%%%%%%%%%%%%%%%%%%%%%%

\begin{frame}{Vectorizing user functions (no side effects)}

% Similar to builtin functions, but with no knowledge of whether arguments need
% packetization. Need to analyze this for each call site.

\end{frame}

%%%%%%%%%%%%%%%%%%%%%%%%%%%%%%%%%%%%%%%%%%%%%%%%%%%%%%%%%%%%%%%%%%%%%%%%%%%%%%%%

\begin{frame}{Vectorizing user functions (with side effects)}

% Need to pass a mask as an extra argument
% Might be simpler to just inline such functions

\end{frame}

%%%%%%%%%%%%%%%%%%%%%%%%%%%%%%%%%%%%%%%%%%%%%%%%%%%%%%%%%%%%%%%%%%%%%%%%%%%%%%%%

\begin{frame}{Interleaved memory optimizations}

\end{frame}

%%%%%%%%%%%%%%%%%%%%%%%%%%%%%%%%%%%%%%%%%%%%%%%%%%%%%%%%%%%%%%%%%%%%%%%%%%%%%%%%

\begin{frame}{SoA to AoS conversion}

\end{frame}




%%%%%%%%%%%%%%%%%%%%%%%%%%%%%%%%%%%%%%%%%%%%%%%%%%%%%%%%%%%%%%%%%%%%%%%%%%%%%%%%

\section*{Conclusion}

\begin{frame}{Summary}

\begin{itemize}
    \item Should be enough to get started and create a functional vectorizer
    \item Many things were not covered in this talk:
    \begin{itemize}
        \item ...
    \end{itemize}
    \item Introduced resources to go further
\end{itemize}

\end{frame}

%%%%%%%%%%%%%%%%%%%%%%%%%%%%%%%%%%%%%%%%%%%%%%%%%%%%%%%%%%%%%%%%%%%%%%%%%%%%%%%%

\begin{frame}{Thank you!}

\begin{itemize}
    \item Q\&A
    \item ...
    \item Happy vectorization! \hspace{1em} <\hspace{0.15em}\penguin, \penguin, \penguin, \penguin, \penguin, \penguin, \penguin, \penguin\hspace{0.1em}>
\end{itemize}

\vspace{5em}

\end{frame}

%%%%%%%%%%%%%%%%%%%%%%%%%%%%%%%%%%%%%%%%%%%%%%%%%%%%%%%%%%%%%%%%%%%%%%%%%%%%%%%%

\end{document}
